\section{Conclusion}
\label{sec:conclusion}

\toolname\ studies a practical question for the agent era: if an individual builder cannot train the model, what can they still do outside the model? Rather than training cheaper models or removing useful context, the project helps repeated context remain cacheable at the harness layer. The current fixed V2 diagnostic supports this narrow claim: after isolating fixture Git state and using robust prompt paths, the cache-friendly condition reduced paid uncached input while preserving validation and task success. The result should still be read conservatively. It is a prefix-cache result, not a universal agent-cost result: tool-output verbosity, extra turns, validation failures, and behavior regressions can still dominate total cost and must be reported separately. The next milestone is therefore not to overstate the prefix result, but to combine it with a separate runtime-output layer and larger paired runs that keep cached tokens, uncached paid input, observed cost, latency, and task success visible side by side.
